\documentclass[10pt,a4paper]{amsart}
\usepackage{amsmath,amssymb,amsthm,graphicx}
% For commutative diagrams we use Xy-pic
\usepackage[all]{xy}

% We use 2cell for 2-commutative diagrams.
\xyoption{2cell}
% Theorem environments.
%
\theoremstyle{plain}
\newtheorem{theorem}[subsection]{Theorem}
\newtheorem{proposition}[subsection]{Proposition}
\newtheorem{lemma}[subsection]{Lemma}

\theoremstyle{definition}
\newtheorem{definition}[subsection]{Definition}
\newtheorem{example}[subsection]{Example}
\newtheorem{exercise}[subsection]{Exercise}
\newtheorem{situation}[subsection]{Situation}

\theoremstyle{remark}
\newtheorem{remark}[subsection]{Remark}
\newtheorem{remarks}[subsection]{Remarks}

\numberwithin{equation}{subsection}

% Macros
%
\def\lim{\mathop{\mathrm{lim}}\nolimits}
\def\colim{\mathop{\mathrm{colim}}\nolimits}
\def\Spec{\mathop{\mathrm{Spec}}}
\def\Hom{\mathop{\mathrm{Hom}}\nolimits}
\def\Ext{\mathop{\mathrm{Ext}}\nolimits}
\def\SheafHom{\mathop{\mathcal{H}\!\mathit{om}}\nolimits}
\def\SheafExt{\mathop{\mathcal{E}\!\mathit{xt}}\nolimits}
\def\Sch{\mathit{Sch}}
\def\Mor{\mathop{\mathrm{Mor}}\nolimits}
\def\Ob{\mathop{\mathrm{Ob}}\nolimits}
\def\Sh{\mathop{\mathit{Sh}}\nolimits}
\def\NL{\mathop{N\!L}\nolimits}
\def\CH{\mathop{\mathrm{CH}}\nolimits}
\def\proetale{{pro\text{-}\acute{e}tale}}
\def\etale{{\acute{e}tale}}
\def\QCoh{\mathit{QCoh}}
\def\Ker{\mathop{\mathrm{Ker}}}
\def\Im{\mathop{\mathrm{Im}}}
\def\Coker{\mathop{\mathrm{Coker}}}
\def\Coim{\mathop{\mathrm{Coim}}}

% Boxtimes
%
\DeclareMathSymbol{\boxtimes}{\mathbin}{AMSa}{"02}

%
% Macros for moduli stacks/spaces
%
\def\QCohstack{\mathcal{QC}\!\mathit{oh}}
\def\Cohstack{\mathcal{C}\!\mathit{oh}}
\def\Spacesstack{\mathcal{S}\!\mathit{paces}}
\def\Quotfunctor{\mathrm{Quot}}
\def\Hilbfunctor{\mathrm{Hilb}}
\def\Curvesstack{\mathcal{C}\!\mathit{urves}}
\def\Polarizedstack{\mathcal{P}\!\mathit{olarized}}
\def\Complexesstack{\mathcal{C}\!\mathit{omplexes}}
% \Pic is the operator that assigns to X its picard group, usage \Pic(X)
% \Picardstack_{X/B} denotes the Picard stack of X over B
% \Picardfunctor_{X/B} denotes the Picard functor of X over B
\def\Pic{\mathop{\mathrm{Pic}}\nolimits}
\def\Picardstack{\mathcal{P}\!\mathit{ic}}
\def\Picardfunctor{\mathrm{Pic}}
\def\Deformationcategory{\mathcal{D}\!\mathit{ef}}
\usepackage{cleveref}
\crefname{originallemma}{original lemma}{original lemmas}
\Crefname{originallemma}{Original lemma}{Original lemmas}
\makeatletter
\newlabel{divisors-lemma-components-locally-finite}{{0BE1}{source}{}{sourcepointer-1}{}}
\newlabel{divisors-lemma-components-locally-finite@cref}{{[originallemma][0][]0BE1}{[1][0][]source}{}{}{}}
\makeatother
\crefname{originallemma}{original lemma}{original lemmas}
\Crefname{originallemma}{Original lemma}{Original lemmas}
\makeatletter
\newlabel{divisors-lemma-divisor-locally-finite}{{02RL}{source}{}{sourcepointer-2}{}}
\newlabel{divisors-lemma-divisor-locally-finite@cref}{{[originallemma][0][]02RL}{[1][0][]source}{}{}{}}
\makeatother
\crefname{originallemma}{original lemma}{original lemmas}
\Crefname{originallemma}{Original lemma}{Original lemmas}
\makeatletter
\newlabel{divisors-lemma-divisor-meromorphic-locally-finite}{{02SG}{source}{}{sourcepointer-3}{}}
\newlabel{divisors-lemma-divisor-meromorphic-locally-finite@cref}{{[originallemma][0][]02SG}{[1][0][]source}{}{}{}}
\makeatother
\crefname{originallemma}{original lemma}{original lemmas}
\Crefname{originallemma}{Original lemma}{Original lemmas}
\makeatletter
\newlabel{lemma-compute-c1}{{02SK}{source}{}{sourcepointer-4}{}}
\newlabel{lemma-compute-c1@cref}{{[originallemma][0][]02SK}{[1][0][]source}{}{}{}}
\makeatother
\crefname{originallemma}{original lemma}{original lemmas}
\Crefname{originallemma}{Original lemma}{Original lemmas}
\makeatletter
\newlabel{lemma-key-formula}{{0AYC}{source}{}{sourcepointer-5}{}}
\newlabel{lemma-key-formula@cref}{{[originallemma][0][]0AYC}{[1][0][]source}{}{}{}}
\makeatother
\crefname{originallemma}{original lemma}{original lemmas}
\Crefname{originallemma}{Original lemma}{Original lemmas}
\makeatletter
\newlabel{section-key}{{0AYB}{source}{}{sourcepointer-6}{}}
\newlabel{section-key@cref}{{[originallemma][0][]0AYB}{[1][0][]source}{}{}{}}
\makeatother
\crefname{originallemma}{original lemma}{original lemmas}
\Crefname{originallemma}{Original lemma}{Original lemmas}
\makeatletter
\newlabel{situation-setup}{{02QL}{source}{}{sourcepointer-7}{}}
\newlabel{situation-setup@cref}{{[originallemma][0][]02QL}{[1][0][]source}{}{}{}}
\makeatother
\title{The Gysin pullback of the divisor of a meromorphic section is the first Chern class of its restriction capped with the fundamental class of the divisor}
\author{The Stacks Project Authors}
\date{stacks-project:chow.tex:local-snapshot-20261012; GNU FDL 1.2 or later}
\begin{document}
\maketitle
\noindent\textit{Editorial note.} All delivered mathematical text is the author's own, copied from the pinned Stacks Project chapter file chow.tex; no statement or proof step was rewritten, reordered or removed.\par\medskip
\noindent\textit{Editorial note.} Preamble provenance: the pinned chapter file opens with \textbackslash{}input\{preamble\} and carries no preamble text of its own, so the project's single Stacks-wide preamble staged at sources/stacks-chow/preamble.tex is excerpted verbatim as two spans, lines 14-18 (Xy-pic setup) and lines 154-217 (the native theorem declarations and mathematical macros). Nothing was added, removed or redefined; the wrapper class and the amsmath/amssymb/amsthm/graphicx base packages are supplied by the assembler.\par\medskip
\noindent\textit{Editorial note.} Same-chapter pointers are located at their \textbackslash{}label lines inside the pinned chapter file; cross-chapter pointers are located in the project's own tag list staged in the snapshot (sources/stacks-chow/tags.txt), at the verbatim lines carrying their canonical tags.\par\medskip
\noindent\textit{Allowed standard prerequisites (editorial).}\begin{itemize}
\item Local finiteness of a divisor and of its components. Divisors, labels divisors-lemma-components-locally-finite, divisors-lemma-divisor-locally-finite and divisors-lemma-divisor-meromorphic-locally-finite, invoked together at proof lines 5370-5372. Lets the divisor div\_N(t) be written as a locally finite sum over integral Z\_i of dimension n - 1 with every irreducible component of D among the Z\_i, which sets up the coefficient computation.
\item The canonical section and the cycle class of a regular section. Divisors, label divisors-definition-regular-meromorphic-ideal-denominators for the canonical section s of O\_X(D) (proof lines 5375-5377); same chapter chow.tex, label lemma-compute-c1, invoked at proof line 5400. Gives s with [D]\_\{n-1\} = sum ord\_\{B\_i\}(f\_i)[Z\_i] and identifies div\_\{L|Z\_i\}(s\_i|Z\_i) with [Z\_i cap D]\_\{n-2\} for those i with Z\_i not contained in D.
\item The key formula for the tame symbol. Same chapter chow.tex, label lemma-key-formula, invoked at proof line 5405; the discussion of the section labelled section-key is used at proof line 5378. Supplies the cycle identity that turns the difference of the two coefficient expressions into a sum of divisors of tame symbols, which is exactly the rational equivalence that finishes the proof.
\end{itemize}\medskip
\ifdefined\hypertarget\hypertarget{sourcepointer-1}{}\fi
\noindent\textit{Source reference.} Original lemma 0BE1: sources/stacks-chow/tags.txt, source lines 14334--14334. Navigation pointer for \textbackslash{}ref\{divisors-lemma-components-locally-finite\} used inside the counted span (line 5370); the argument that consumes it is retained in the span.\par\medskip
\ifdefined\hypertarget\hypertarget{sourcepointer-2}{}\fi
\noindent\textit{Source reference.} Original lemma 02RL: sources/stacks-chow/tags.txt, source lines 3628--3628. Navigation pointer for \textbackslash{}ref\{divisors-lemma-divisor-locally-finite\} used inside the counted span (line 5371); the argument that consumes it is retained in the span.\par\medskip
\ifdefined\hypertarget\hypertarget{sourcepointer-3}{}\fi
\noindent\textit{Source reference.} Original lemma 02SG: sources/stacks-chow/tags.txt, source lines 3659--3659. Navigation pointer for \textbackslash{}ref\{divisors-lemma-divisor-meromorphic-locally-finite\} used inside the counted span (line 5372); the argument that consumes it is retained in the span.\par\medskip
\ifdefined\hypertarget\hypertarget{sourcepointer-4}{}\fi
\noindent\textit{Source reference.} Original lemma 02SK: sources/stacks-chow/pinned-original-source.tex, source lines 4198--4198. Navigation pointer for \textbackslash{}ref\{lemma-compute-c1\} used inside the counted span (line 5400); the argument that consumes it is retained in the span.\par\medskip
\ifdefined\hypertarget\hypertarget{sourcepointer-5}{}\fi
\noindent\textit{Source reference.} Original lemma 0AYC: sources/stacks-chow/pinned-original-source.tex, source lines 4746--4746. Navigation pointer for \textbackslash{}ref\{lemma-key-formula\} used inside the counted span (line 5405); the argument that consumes it is retained in the span.\par\medskip
\ifdefined\hypertarget\hypertarget{sourcepointer-6}{}\fi
\noindent\textit{Source reference.} Original lemma 0AYB: sources/stacks-chow/pinned-original-source.tex, source lines 4686--4686. Navigation pointer for \textbackslash{}ref\{section-key\} used inside the counted span (line 5378); the argument that consumes it is retained in the span.\par\medskip
\ifdefined\hypertarget\hypertarget{sourcepointer-7}{}\fi
\noindent\textit{Source reference.} Original lemma 02QL: sources/stacks-chow/pinned-original-source.tex, source lines 1372--1372. Navigation pointer for \textbackslash{}ref\{situation-setup\} used inside the counted span (line 5353); the argument that consumes it is retained in the span.\par\medskip
\begin{lemma}
\label{lemma-gysin-factors-general}
Let $(S, \delta)$ be as in Situation \ref{situation-setup}.
Let $X$ be locally of finite type over $S$.
Let $X$ be integral and $n = \dim_\delta(X)$.
Let $i : D \to X$ be an effective Cartier divisor.
Let $\mathcal{N}$ be an invertible $\mathcal{O}_X$-module
and let $t$ be a nonzero meromorphic section of $\mathcal{N}$.
Then $i^*\text{div}_\mathcal{N}(t) = c_1(\mathcal{N}|_D) \cap [D]_{n - 1}$
in $\CH_{n - 2}(D)$.
\end{lemma}

\begin{proof}
Write $\text{div}_\mathcal{N}(t) = \sum \text{ord}_{Z_i, \mathcal{N}}(t)[Z_i]$
for some integral closed subschemes $Z_i \subset X$ of $\delta$-dimension
$n - 1$. We may assume that the family $\{Z_i\}$ is locally
finite, that $t \in \Gamma(U, \mathcal{N}|_U)$ is a generator
where $U = X \setminus \bigcup Z_i$, and that every irreducible component
of $D$ is one of the $Z_i$, see
Divisors, Lemmas \ref{divisors-lemma-components-locally-finite},
\ref{divisors-lemma-divisor-locally-finite}, and
\ref{divisors-lemma-divisor-meromorphic-locally-finite}.

\medskip\noindent
Set $\mathcal{L} = \mathcal{O}_X(D)$. Denote
$s \in \Gamma(X, \mathcal{O}_X(D)) = \Gamma(X, \mathcal{L})$
the canonical section. We will apply the discussion of
Section \ref{section-key} to our current situation.
For each $i$ let $\xi_i \in Z_i$ be its generic point. Let
$B_i = \mathcal{O}_{X, \xi_i}$. For each $i$ we pick generators
$s_i \in \mathcal{L}_{\xi_i}$ and $t_i \in \mathcal{N}_{\xi_i}$
over $B_i$ but we insist that we pick $s_i = s$ if $Z_i \not \subset D$.
Write $s = f_i s_i$ and $t = g_i t_i$ with $f_i, g_i \in B_i$.
Then $\text{ord}_{Z_i, \mathcal{N}}(t) = \text{ord}_{B_i}(g_i)$.
On the other hand, we have $f_i \in B_i$ and
$$
[D]_{n - 1} = \sum \text{ord}_{B_i}(f_i)[Z_i]
$$
because of our choices of $s_i$. We claim that
$$
i^*\text{div}_\mathcal{N}(t) =
\sum \text{ord}_{B_i}(g_i) \text{div}_{\mathcal{L}|_{Z_i}}(s_i|_{Z_i})
$$
as cycles. More precisely, the right hand side is a cycle
representing the left hand side. Namely, this is clear by our
formula for $\text{div}_\mathcal{N}(t)$ and the fact that
$\text{div}_{\mathcal{L}|_{Z_i}}(s_i|_{Z_i}) = [Z(s_i|_{Z_i})]_{n - 2} =
[Z_i \cap D]_{n - 2}$ when $Z_i \not \subset D$ because in
that case $s_i|_{Z_i} = s|_{Z_i}$ is a regular section, see
Lemma \ref{lemma-compute-c1}. Similarly,
$$
c_1(\mathcal{N}) \cap [D]_{n - 1} =
\sum \text{ord}_{B_i}(f_i) \text{div}_{\mathcal{N}|_{Z_i}}(t_i|_{Z_i})
$$
The key formula (Lemma \ref{lemma-key-formula}) gives the equality
$$
\sum \left(
\text{ord}_{B_i}(f_i) \text{div}_{\mathcal{N}|_{Z_i}}(t_i|_{Z_i}) -
\text{ord}_{B_i}(g_i) \text{div}_{\mathcal{L}|_{Z_i}}(s_i|_{Z_i}) \right) =
\sum \text{div}_{Z_i}(\partial_{B_i}(f_i, g_i))
$$
of cycles. If $Z_i \not \subset D$, then $f_i = 1$ and hence
$\text{div}_{Z_i}(\partial_{B_i}(f_i, g_i)) = 0$. Thus we get a rational
equivalence between our specific cycles representing
$i^*\text{div}_\mathcal{N}(t)$ and $c_1(\mathcal{N}) \cap [D]_{n - 1}$
on $D$. This finishes the proof.
\end{proof}

\end{document}
